% !TeX root = ../main.tex

\chapter{Introduction}\label{chapter:introduction}

\section{Motivation}

Photorealistic capture of real-world environments has progressed rapidly in
the last few years. \ac{NeRF}-based methods first demonstrated that a scene
can be reconstructed implicitly from a handful of posed photographs, and
\ac{3DGS}~\parencite{kerbl3Dgaussians} has since become the dominant
practical representation for high-fidelity novel view synthesis: rendering
is real-time, training takes minutes instead of hours, and the output is an
explicit point cloud of anisotropic Gaussians that can be streamed to a
browser.

A natural next question is whether scenes captured this way can also be
\emph{edited}. Concretely: given a 3DGS reconstruction of a real room, can
we insert a new mesh-based object --- a chair, a lamp, an appliance --- so
that the result looks like it was always there? This use case is central to
interior design, e-commerce visualization, and game-asset content creation,
but the implicit, per-Gaussian representation does not expose the surface
geometry, surface materials, or environment lighting that classical
compositing relies on.

\section{Problem statement}

This thesis investigates one concrete instance of the problem: insertion of
a known textured mesh into a known 3DGS room, such that
\begin{enumerate}
  \item shading on the inserted object is consistent with the room's
        captured illumination,
  \item soft contact shadows are produced on the floor,
  \item rendering remains compatible with the existing real-time 3DGS
        viewers and does not require retraining the host scene.
\end{enumerate}

\section{Contributions}

The main contributions of this work, packaged as the \emph{InsertGS}
pipeline, are:

\begin{itemize}
  \item A GPU rasterizer (\texttt{gs\_to\_cubemap\_gsplat.py}) that converts
        a 3DGS room into a six-face cubemap using proper \ac{EWA}
        splatting, suitable for \ac{IBL}.
  \item A \ac{PBR} rendering stage in Blender that consumes the cubemap as
        an environment map and produces shaded RGB plus alpha and shadow
        passes for the inserted mesh.
  \item A compositing module that places the rendered object back into
        novel 3DGS viewpoints while preserving depth ordering with the
        surrounding splats.
  \item A unified Python entry point (\texttt{main.py}) exposing the
        pipeline both as a \ac{CLI} for batch experiments and as a
        \ac{GUI} dashboard for interactive use, plus a lightweight
        SuperSplat-based web viewer for inspecting \ac{SPZ} and \ac{PLY}
        scenes in the browser.
\end{itemize}

\section{Outline}

\Cref{chapter:related} reviews 3DGS, NeRF-based editing, and prior work on
inserting objects into captured radiance fields. \Cref{chapter:method}
develops the pipeline in three stages: cubemap extraction, PBR rendering,
and compositing. \Cref{chapter:implementation} describes the engineering
artefacts: the case-based directory layout, the Gradio dashboard, and the
SPZ web viewer. \Cref{chapter:results} reports qualitative and quantitative
results on a curated set of indoor scenes, and
\Cref{chapter:conclusion} summarizes findings and limitations.
